% Learning InitechBase, Chapter 1: Quick Start (PUB-REF-0001 AN-16; annex
% Sec 10).  Every panel replayed (examples/l1-*.dot).
\chapter{Quick Start}
\Gates{test-pubs-examples, test-samir-repl, test-samir-query}
In this lesson you open the staff database of Initech Systems Corporation,
look at its records, ask InitechBase two questions, and leave InitechBase.
You use five commands: \cmd{USE}, \cmd{LIST}, \cmd{DISPLAY}, \cmd{?} and
\cmd{QUIT}.

\begin{ChapterCovers}
\item The dot prompt
\item Opening a database with \cmd{USE}
\item Listing records with \cmd{LIST}
\item Displaying one record with \cmd{DISPLAY}
\item Asking questions with \cmd{?}
\item Leaving InitechBase with \cmd{QUIT}
\end{ChapterCovers}

\ATEHead{Preparing for This Chapter}
Copy the files of the Sample Programs and Utilities diskette to the disk you
work on, as \doctitle{Getting Started} describes, and start InitechBase from
that disk.  The screen shows the dot prompt: a period and a blank, with the
cursor after them.

\ATEHead{The Dot Prompt}
The dot prompt means that InitechBase is waiting for a command.  You type a
command, press \key{Enter}, and InitechBase carries it out and shows the dot
prompt again.  Commands are words of English: \cmd{USE} a database,
\cmd{LIST} its records, \cmd{QUIT}.

If you make a mistake while typing, press \key{Backspace} to erase it.  If you
press \key{Enter} on a command InitechBase does not know, it says so; nothing
is changed, and you can type the command again.

\ATEHead{Opening a Database}
A database is a file of records.  Each record holds the information about one
thing -- here, one employee of Initech -- in fields: the employee number, the
last name, the first name, the department, the telephone extension, the date
of hire, the salary, whether the employee is exempt from overtime, and a note.
Before you can look at a database, you open it with \cmd{USE}.
\begin{Steps}
\item Type \type{USE STAFF} and press \key{Enter}.
\ATEResponse{InitechBase opens \file{STAFF.DBF} and shows the dot prompt again.}
\end{Steps}
\DotExample{examples/l1-use}
\cmd{USE} added the extension \file{.DBF} to the name for you.  The database
stays open until you close it or leave InitechBase.

\ATEHead{Listing the Records}
\begin{Steps}[resume]
\item Type \type{LIST LASTNAME, FIRSTNAME, EXT} and press \key{Enter}.
\ATEResponse{InitechBase lists the names and telephone extensions of all the
employees.}
\end{Steps}
\DotExample{examples/l1-list}
The first line is the heading: \texttt{Record\#} and the names of the fields
you asked for.  Each record follows on its own line.  The number at the left
is the record number, the position of the record in the database.

The asterisk in front of record 13 means that the record is marked for
deletion.  Chapter~5 shows you how records are marked and how the marked
records are removed.

\begin{AteNote}
\cmd{LIST} with no field names lists every field of every record.  The records
of \file{STAFF.DBF} are wider than the screen, so each takes two lines.  Name
the fields you want, as you did in step 2.
\end{AteNote}

\ATEHead{Displaying One Record}
InitechBase keeps track of one record of the open database, the current
record.  \cmd{DISPLAY} shows it.  \cmd{GO} makes another record current.
\begin{Steps}[resume]
\item Type \type{GO 6} and press \key{Enter}.
\item Type \type{DISPLAY} and press \key{Enter}.
\ATEResponse{InitechBase displays every field of record 6.}
\end{Steps}
\DotExample{examples/l1-display}
The record of Tom Smykowski of Customer Relations is too wide for the
screen; the end of his note continues on the next line.

\ATEHead{Asking a Question}
\cmd{?} displays the value of a field, or of a calculation, without the
heading and the record number.
\begin{Steps}[resume]
\item Type \type{GO 5} and press \key{Enter}.
\item Type \type{? FIRSTNAME, LASTNAME} and press \key{Enter}.
\item Type \type{? SALARY / 12} and press \key{Enter}.
\ATEResponse{InitechBase displays the monthly salary of Milton Waddams.}
\end{Steps}
\DotExample{examples/l1-question}
\cmd{?} displays its answer and leaves the cursor at the end of it, so the next
dot prompt follows on the same line.

\ATEHead{Choosing Records}
A \kw{FOR} clause makes \cmd{LIST} list only the records you describe.
\begin{Steps}[resume]
\item Type \type{LIST LASTNAME, DEPT, EXT FOR DEPT = 'CUST'} and press
\key{Enter}.
\ATEResponse{InitechBase lists the employees of Customer Relations.}
\end{Steps}
\DotExample{examples/l1-for}
The department code is in quotation marks because it is a character string,
not the name of a field.  Chapter~2 tells you more about conditions.

\ATEHead{Leaving InitechBase}
\begin{Steps}[resume]
\item Type \type{QUIT} and press \key{Enter}.
\ATEResponse{InitechBase closes the database and returns to InitechDOS.}
\end{Steps}
\DotExample{examples/l1-quit}
Always leave InitechBase with \cmd{QUIT} before you turn off the computer:
\cmd{QUIT} closes the files properly.

\ATEHead{Where Do You Go From Here?}
You have opened a database, listed its records, displayed one of them, asked
two questions and left InitechBase.  In Chapter~2 you choose the fields and
the records that \cmd{LIST} shows.

\ATEHead{Summary}
\begin{enumerate}
\item \cmd{USE} \xp{database} opens a database.
\item \cmd{LIST} \xp{field list} lists the fields of every record.
\item \cmd{GO} \xp{n} makes record \xp{n} the current record; \cmd{DISPLAY} shows
  it.
\item \cmd{?} \xp{expression list} displays values.
\item \kw{FOR} \xp{condition} chooses the records a command works on.
\item \cmd{QUIT} closes every file and leaves InitechBase.
\end{enumerate}
